\scanpage{118}
% line head | text | 67
\mspage{67}
% line 1–3 | text | qui soit symétrique en U, T
qui soit symétrique en $U,T$, j'ai envie
de regarder d'abord la situation ``mixte'', d'un
diagramme en traits pleins
% line 4–6 | diagram | (α) U to X from T over U prime to X prime from T prime
\diagram{\begin{tikzcd}[ampersand replacement=\&,column sep=small]
\& U\arrow[r,hook,"i"]\arrow[d,"f_U"'] \& X\arrow[d,dashed,"f"] \& T\arrow[l,hook',"j"']\arrow[d,"f_T"]\\
{(\unsure{\alpha|a})} \& U'\arrow[r,dashed,hook,"i'"] \& X' \& T'\arrow[l,dashed,hook',"j'"'] {,}
\end{tikzcd}}
% line 7–14 | text | qu'on aimerait compléter
qu'on aimerait compléter ``canoniquement'', si faire
se peut, en un morphisme de triples (i.e. en
une flèche de $\unsure{\mathscr T|\mathcal T}$). L'idée naturelle,
c'est que le $X'$ qui s'y \unsure{insère} \unsure{devrait|serait}
la \unsure{limite} \unsure{inductive}, dans $\mathrm{Cat}$, du diagramme
initial. Si \inserted{$T'=T$,} $f_T=\mathrm{id}_T$, cela signifie donc
que $X'$ s'identifie à la somme amalgamée
$X\amalg_{U\unsure{\alpha|{,}\alpha}}U'$, et inversement si $U'=U$, $f_U=\mathrm{id}_U$.

% line 15–18 | text | Désignons par X prime
Désignons par $X'$ la somme amalgamée et
\deleted{\unsure{défin.}} explicitons, pour $X''$ quelconque \unsure{de|dans} $\mathrm{Cat}$,
\inserted{l'ens} des flèches \unsure{?} de $X'$ dans $X''$, \deleted{\unsure{celui-ci}} comme
foncteur en $X''$. C'est l'ens des triples
% line 19–21 left | diagram | U and T over U prime and T prime, X mapping to X double prime
\diagram{\begin{tikzcd}[ampersand replacement=\&,column sep=small]
 U\arrow[r,hook,"i"]\arrow[d,"f_U"'] \& X\arrow[dd,"g_X"'] \& \deleted{\unsure{U}}T\arrow[l,hook',"j"']\arrow[d,"f_T"]\\
 U'\arrow[dr,"g_{U'}"'] \& \& T'\arrow[dl,"g_{T'}"]\\
 \& X'' \&
\end{tikzcd}}
% line 19–21 right | display | (*) triple g_X, g_U prime, g_T prime and equations
\[
\begin{aligned}
&(g_X,g_{U'},g_{T'})\quad
\left|\begin{aligned}
&g_X:X\to X''\\
&g_{U'}:U'\to X''\\
&g_{T'}:T'\to X''
\end{aligned}\right.\qquad\text{avec}\quad
\begin{aligned}
g_X\circ i&=g_{U'}f_U\\
g_X\circ j&=g_{T'}f_T
\end{aligned}
\end{aligned}\tag*{(*)}
\]
% line 22–25 | text | Notons que pour un tel triple
Notons que pour un tel triple,
$g_X$ est partiellement connu \inserted{déjà} en termes
de $g_{U'},g_{T'}$, par ses restrictions \inserted{$g_X\circ i$ et $g_X\circ j$} à
$U$ et $T$ respectivement, par les formules
% line 26 | display | g_U, g_T
\[
g_U=g_{U'}f_U,\qquad g_T=g_{T'}f_T
\]
% archive sheet number: 118
